% ==========================================================================
%  Chapter 73 -- Arc-Length Control for Nonlinear Analysis
% ==========================================================================

\chapter{Arc-Length Control for Nonlinear Analysis}
\label{ch:arc-length}


% =========================================================================
\section{Introduction}
\label{sec:arc-intro}

Load-controlled and displacement-controlled Newton–Raphson methods
fail at limit points (snap-through) and bifurcation points
(snap-back) on the equilibrium path.  At a limit point the tangent
stiffness matrix becomes singular in the loading direction; at a
snap-back the displacement path reverses, requiring negative load
increments that standard monotonic schemes cannot produce.

The arc-length method---originally proposed by Riks (1979) and refined
by Crisfield (1981, 1983)---augments the equilibrium system with a
\emph{constraint equation} that controls the step length along the
equilibrium path in the combined load--displacement space.  This
chapter documents the mathematical formulation and its implementation
in the \texttt{ArcLengthControl} and \texttt{ArcLengthSolver} classes.


% =========================================================================
\section{Augmented System}
\label{sec:arc-augmented}

We seek the equilibrium of a discretised structure under a
proportional reference load $\hat{\mathbf{f}}$ scaled by a single
load parameter $\lambda$:
\begin{equation}
  \label{eq:arc-equilibrium}
  \mathbf{R}(\mathbf{u}, \lambda)
  = \mathbf{f}_{\text{int}}(\mathbf{u})
    - \lambda \, \hat{\mathbf{f}} = \mathbf{0}.
\end{equation}
The unknowns are the displacement vector
$\mathbf{u} \in \mathbb{R}^n$ and the scalar load parameter
$\lambda \in \mathbb{R}$, totalling $n + 1$ equations for $n + 1$
unknowns once a constraint is appended.

% -------------------------------------------------------------------------
\subsection{Constraint Equations}

Two variants are implemented:

\paragraph{Cylindrical.}
The cylindrical arc-length constraint (Crisfield, 1981) enforces a
constant displacement-space step:
\begin{equation}
  \label{eq:arc-cylindrical}
  g(\Delta\mathbf{u}, \Delta\lambda)
  = \Delta\mathbf{u}^T \Delta\mathbf{u}
    + \psi^2 \, \Delta\lambda^2 \,
      \hat{\mathbf{f}}^T \hat{\mathbf{f}}
  - \Delta\ell^2 = 0,
\end{equation}
where $\Delta\ell$ is the prescribed arc-length radius and $\psi$ is a
scaling factor that controls the relative weight of the load component
($\psi = 0$ reduces to a purely displacement-based norm,
$\psi = 1$ gives the ``spherical'' version in the original coordinates).

\paragraph{Spherical.}
The spherical variant (Crisfield, 1981) uses the full
$(n+1)$-dimensional norm:
\begin{equation}
  \label{eq:arc-spherical}
  g = \Delta\mathbf{u}^T \Delta\mathbf{u}
    + \psi^2 \, \Delta\lambda^2 \,
      \hat{\mathbf{f}}^T \hat{\mathbf{f}}
  - \Delta\ell^2 = 0.
\end{equation}
Both constraints share the same algebraic form but differ in the
root-selection strategy of the corrector (see below).


% =========================================================================
\section{Predictor Step (Tangent Predictor)}
\label{sec:arc-predictor}

At the beginning of each new step $k$, the tangent stiffness at the
last converged state $(\mathbf{u}_k, \lambda_k)$ is factored and the
load-sensitivity vector is computed:
\begin{equation}
  \label{eq:arc-duf}
  \mathbf{K}_t \, \delta\mathbf{u}_f = \hat{\mathbf{f}}.
\end{equation}

The predictor increment satisfies the constraint exactly:
\begin{equation}
  \label{eq:arc-pred-dl}
  \Delta\lambda_{\text{pred}} = \pm \frac{\Delta\ell}
    {\sqrt{\delta\mathbf{u}_f^T \delta\mathbf{u}_f
           + \psi^2 \,
             \hat{\mathbf{f}}^T \hat{\mathbf{f}}}}.
\end{equation}

The sign is chosen to maintain a consistent traversal direction:
\begin{equation}
  \label{eq:arc-pred-sign}
  \operatorname{sign}(\Delta\lambda_{\text{pred}})
  = \operatorname{sign}(\Delta\lambda_{k-1}),
\end{equation}
where $\Delta\lambda_{k-1}$ is the \emph{converged} load increment from
the previous step.  At the very first step, a positive sign (loading)
is assumed.  After the predictor:
\begin{equation}
  \Delta\mathbf{u}_{\text{pred}} = \Delta\lambda_{\text{pred}} \,
  \delta\mathbf{u}_f.
\end{equation}


% =========================================================================
\section{Corrector Iterations (Sherman--Morrison Decomposition)}
\label{sec:arc-corrector}

At each Newton corrector iteration $j$ within step $k$, the current
state is $(\mathbf{u}_k^{(j)}, \lambda_k^{(j)})$ with accumulated
increments $(\Delta\mathbf{u}^{(j)}, \Delta\lambda^{(j)})$.  The
residual is:
\begin{equation}
  \mathbf{R}^{(j)} = \mathbf{f}_{\text{int}}(\mathbf{u}_k^{(j)})
  - \lambda_k^{(j)} \, \hat{\mathbf{f}}.
\end{equation}

Two right-hand-side solves are performed:
\begin{equation}
  \label{eq:arc-two-solves}
  \mathbf{K}_t^{(j)} \, \delta\mathbf{u}_R = -\mathbf{R}^{(j)},
  \qquad
  \mathbf{K}_t^{(j)} \, \delta\mathbf{u}_f = \hat{\mathbf{f}}.
\end{equation}
The total correction is decomposed as:
\begin{equation}
  \label{eq:arc-du-decomp}
  \delta\mathbf{u} = \delta\mathbf{u}_R
  + \delta\lambda \, \delta\mathbf{u}_f,
\end{equation}
where $\delta\lambda$ is determined from the constraint linearisation.

% -------------------------------------------------------------------------
\subsection{Cylindrical Corrector}

Linearising \eqref{eq:arc-cylindrical} about
$(\Delta\mathbf{u}^{(j)}, \Delta\lambda^{(j)})$ and substituting
\eqref{eq:arc-du-decomp}:
\begin{equation}
  \label{eq:arc-cyl-dl}
  \delta\lambda = -\frac{
    \Delta\mathbf{u}^{(j)\,T} \delta\mathbf{u}_R
  }{
    \Delta\mathbf{u}^{(j)\,T} \delta\mathbf{u}_f
    + \psi^2 \, \Delta\lambda^{(j)} \,
      \hat{\mathbf{f}}^T \hat{\mathbf{f}}
  }.
\end{equation}
This is the Sherman--Morrison formula for the bordered system.

% -------------------------------------------------------------------------
\subsection{Spherical Corrector}

The full constraint \eqref{eq:arc-spherical} yields a quadratic in
$\delta\lambda$:
\begin{equation}
  \label{eq:arc-sph-quad}
  a\,\delta\lambda^2 + b\,\delta\lambda + c = 0,
\end{equation}
with:
\begin{align}
  a &= \delta\mathbf{u}_f^T \delta\mathbf{u}_f
       + \psi^2 \, \|\hat{\mathbf{f}}\|^2, \\
  b &= 2 \bigl(
         (\Delta\mathbf{u}^{(j)} + \delta\mathbf{u}_R)^T
         \delta\mathbf{u}_f
         + \psi^2 \, \Delta\lambda^{(j)} \,
           \|\hat{\mathbf{f}}\|^2
       \bigr), \\
  c &= (\Delta\mathbf{u}^{(j)} + \delta\mathbf{u}_R)^T
       (\Delta\mathbf{u}^{(j)} + \delta\mathbf{u}_R)
       + \psi^2 \, \Delta\lambda^{(j)\,2} \,
         \|\hat{\mathbf{f}}\|^2
       - \Delta\ell^2.
\end{align}
The root with the smaller absolute value is chosen:
\begin{equation}
  \delta\lambda = \frac{-b \pm \sqrt{b^2 - 4ac}}{2a},
  \qquad \text{selecting } |\delta\lambda| \to \min.
\end{equation}
When the discriminant is negative (the corrector cannot reach the
arc-length sphere), the cylindrical formula
\eqref{eq:arc-cyl-dl} is used as a fallback.


% =========================================================================
\section{Convergence Check and State Commit}
\label{sec:arc-converge}

At each corrector iteration:
\begin{equation}
  \Delta\mathbf{u}^{(j+1)} = \Delta\mathbf{u}^{(j)}
    + \delta\mathbf{u}_R + \delta\lambda \, \delta\mathbf{u}_f,
  \qquad
  \Delta\lambda^{(j+1)} = \Delta\lambda^{(j)} + \delta\lambda.
\end{equation}

Convergence is declared when the residual norm falls below the
tolerance:
\begin{equation}
  \label{eq:arc-tol}
  \frac{\|\mathbf{R}^{(j+1)}\|}{\|\hat{\mathbf{f}}\|} < \epsilon_r.
\end{equation}

Upon convergence the material state is committed and the step result
$(\lambda, \Delta\lambda, \Delta\ell)$ is stored.  On divergence the
material state is reverted and the step is reported as failed.


% =========================================================================
\section{Implementation: \texttt{ArcLengthControl}}
\label{sec:arc-control-impl}

The \texttt{ArcLengthControl} struct in \texttt{IncrementalControl.hh}
satisfies the existing \texttt{IncrementalControlPolicy} concept:
\begin{equation*}
  \texttt{scheme.apply}(p, \hat{\mathbf{f}},
  \mathbf{f}_{\text{ext}}, \text{model}),
\end{equation*}
where $p \in [0, 1]$ is the load proportionality parameter.
Internally, the control tracks:
\begin{itemize}
  \item $\Delta\ell$: arc-length radius,
  \item $\psi$: scaling parameter,
  \item $\lambda$: current total load factor,
  \item $\Delta\mathbf{u}$: accumulated step displacement (PETSc Vec),
  \item $\Delta\lambda_{\text{prev}}$: last converged load increment
        (for sign determination),
  \item $\|\hat{\mathbf{f}}\|^2$: cached reference load norm.
\end{itemize}

The control manages a PETSc \texttt{Vec} for $\Delta\mathbf{u}$ with
proper lifecycle (destructor, move semantics, deleted copy).

\paragraph{Variant enum.}
\texttt{ArcLengthVariant::Cylindrical} and
\texttt{ArcLengthVariant::Spherical} select the constraint equation
used in the corrector.


% =========================================================================
\section{Implementation: \texttt{ArcLengthSolver}}
\label{sec:arc-solver-impl}

The standalone \texttt{ArcLengthSolver} provides a complete
SNES-free solver for tracing equilibrium paths.  It manages its own
PETSc vectors ($\mathbf{U}$, $\mathbf{R}$, $\hat{\mathbf{f}}$,
$\mathbf{f}_{\text{ext}}$), matrix ($\mathbf{K}$), and KSP (direct
LU via PCLU).

\texttt{solve\_step()} implements the full algorithm:
\begin{enumerate}
  \item Assemble tangent $\mathbf{K}_t$ at current state.
  \item Solve $\mathbf{K}_t \, \delta\mathbf{u}_f = \hat{\mathbf{f}}$.
  \item Tangent predictor: compute $\Delta\lambda_{\text{pred}}$
    via \eqref{eq:arc-pred-dl}, update
    $\Delta\mathbf{u}_{\text{pred}}$.
  \item Corrector loop (up to 50 iterations):
    \begin{enumerate}
      \item Assemble residual $\mathbf{R}$ and tangent $\mathbf{K}_t$.
      \item Solve \eqref{eq:arc-two-solves} for
        $\delta\mathbf{u}_R$ and $\delta\mathbf{u}_f$.
      \item Compute $\delta\lambda$ from constraint
        (cylindrical or spherical).
      \item Update $\Delta\mathbf{u}$, $\Delta\lambda$,
        $\mathbf{U}$, $\lambda$.
      \item Check convergence \eqref{eq:arc-tol}.
    \end{enumerate}
  \item On convergence: commit material state.
  \item On divergence: revert material state, report failure.
\end{enumerate}

The result struct \texttt{ArcLengthResult} reports: convergence flag,
$\lambda$, $\Delta\lambda$, $\Delta\ell$, iteration count, and
final residual norm.

\texttt{solve\_n\_steps(int n)} provides a convenience wrapper that
traces $n$ steps sequentially with progress reporting.


% =========================================================================
\section{Summary}
\label{sec:arc-summary}

The arc-length implementation provides two complementary components:
\begin{enumerate}
  \item \texttt{ArcLengthControl}: a lightweight constraint that plugs
    into the existing \texttt{IncrementalControlPolicy} framework for
    use within existing solvers.
  \item \texttt{ArcLengthSolver}: a self-contained solver with its own
    PETSc infrastructure for standalone equilibrium-path tracing.
\end{enumerate}

Both support cylindrical and spherical constraint variants, the
Sherman--Morrison decomposition for the bordered system, and sign
tracking for consistent traversal past limit and turning points.  The
implementation is general enough to work with any \texttt{Model<>}
configuration, including the mixed concrete/rebar models introduced
in Chapter~\ref{ch:mixed-model-rebar}.
